\documentclass[11pt]{article}
\usepackage[margin=25mm]{geometry}
\usepackage{fontspec}
\setmainfont{Latin Modern Roman}
\usepackage{amsmath,amssymb,mathtools}
\usepackage{unicode-math}
\setmathfont{Latin Modern Math}
\usepackage{xurl,hyperref}
\hypersetup{colorlinks=true,urlcolor=blue}
\setcounter{secnumdepth}{0}
\setlength{\parindent}{0pt}
\setlength{\parskip}{5pt}
\setlength{\emergencystretch}{3em}
\newcommand{\sref}[2]{\hyperlink{source:#1:#2}{[S#2]}}
\newcommand{\cataloguescope}[1]{\paragraph{Recorded target.}#1}
\newcommand{\R}{\mathbb R}
\newcommand{\C}{\mathbb C}
\newcommand{\N}{\mathbb N}
\newcommand{\Z}{\mathbb Z}
\newcommand{\Q}{\mathbb Q}
\newcommand{\F}{\mathbb F}
\newcommand{\eps}{\varepsilon}
\begin{document}
\section{Continuation of Nordheim Condensation Beyond Dirac-Mass Formation}\label{statement}

\subsection{Definitions and mathematical statement}

\paragraph{Shared notation.}$\mathbb N=\{1,2,\ldots\}$; the standard number fields have their usual meanings. All additional conventions are stated locally.


The Nordheim equation describes a spatially homogeneous Bose gas through an energy distribution. The energy variable is $\omega\geq0$. Condensation corresponds to a measure with an atom at zero energy; $n(t)$ is its mass. In the source normalization the particle measure uses the weight $\sqrt\omega\,d\omega$, so the expression $\delta_0/\sqrt\omega$ is a symbolic way to describe that atom, not ordinary multiplication of distributions.

\paragraph{Catalogue wording retained for traceability.}
After finite-time formation of a Dirac mass at $\omega=0$ in the Nordheim equation, determine whether the solution can be continued as $f(t,\omega)=n(t)\delta_0/\sqrt\omega+h(t,\omega)$ with $h(t,\cdot)$ a function.

\paragraph{Scope and source qualification.}A measure-valued weak formulation is necessary at the atom. Conservation laws, positivity and any uniqueness claim must be justified in that formulation.
\subsection{Short English statement}
Determine whether the Nordheim evolution can continue after condensation as a zero-energy atom plus a regular energy-density component.
\subsection{Research attempt}
\paragraph{Model, date and investigated scope.}
OpenAI GPT-6 Astra (configured model \texttt{gpt-6-astra}), ultra reasoning; generated 11 October 2026. Constructed positive atom-plus-density measures with fixed mass and energy to demonstrate that conservation cannot select a post-condensation evolution.

\paragraph{Sources and literature check.}
Read the full supplied definition, catalogue wording, source qualifications and listed references for this question. No fresh status or novelty search was performed. The argument below proves its auxiliary statements directly and does not import an unverified external theorem.

\paragraph{Testing continuation through conserved moments.} Work with the particle measure \(G=n\delta_0+\sqrt\omega h(\omega)\,d\omega\), rather than multiplying the atom by an undefined reciprocal square root. Fix particle number \(M>0\) and energy \(E>0\). For any \(0\leq n<M\), set \(a=(M-n)/E\) and
\[
 g_n(\omega)=(M-n)a e^{-a\omega},\qquad
 G_n=n\delta_0+g_n(\omega)\,d\omega.
\]
Integration gives \(\int dG_n=M\) and \(\int\omega\,dG_n=E\). The corresponding regular occupation is \(h_n=g_n/\sqrt\omega\) on \(\omega>0\); its singularity at zero is locally integrable even without the particle weight. Thus conservation, positivity and an atom-plus-function decomposition allow every value of the atom mass below \(M\).
\paragraph{Failure of moment closure.} Any differentiable path \(n(t)\in[0,M)\) inserted in this formula preserves these two moments. Paths with increasing, decreasing or constant condensate mass all pass this test. They are not asserted to solve the Nordheim equation: the calculation disproves the attempted inference of a condensate evolution law from conserved moments alone. At \(E=0\), by contrast, a positive measure supported on \([0,\infty)\) with zero first moment is supported at zero, so \(G=M\delta_0\).
\paragraph{Remaining source task.} The first obstacle is defining and controlling the collision functional on a measure containing the zero-energy atom, including atom--density interactions and their flux into the atom. None of the constructed paths verifies that weak equation. Existence after condensation, persistence of an absolutely continuous remainder, positivity and any uniqueness assertion in the source formulation remain open within this attempt. The source qualification in [S2] is essential to this distinction.

\paragraph{Verification and adversarial audit.}
Checked both moments explicitly, local integrability of h, n<M and the distinct zero-energy boundary case. Model self-audit only; no independent human mathematical verification. Applied contribution prompt A before drafting and B afterwards. Checked both moments explicitly, local integrability of h, n<M and the distinct zero-energy boundary case. First remaining gap: Define the Nordheim collision functional at the atom and control its atom--density flux in an actual weak continuation. The inherited source status is preserved; original-target completion remains zero.

Coordinator checked the condensate-plus-exponential moments for every allowed atom mass and the zero-energy boundary case. This additional check was model-assisted, not independent human verification.

\paragraph{Tools, process and contributions.}
Python 3 and local file tools used to read definitions and serialize manuscripts. Calculations are symbolic reasoning; no numerical experiment or theorem prover. Contribution prompt A was followed by prompt B. GPT-6 Astra with ultra reasoning drafted the research and performed its stated self-audit; any additional coordinator contribution is identified in the verification and audit. A GPT-6-based Codex coordinator assembled the manuscript and checked repository integration; its exact variant and reasoning configuration were not exposed. No human mathematical contribution, independent human proof check or formal proof certification is supplied. Complete derivations and the consulted source references are included here; private reasoning transcripts are not published.

\paragraph{Final status and first remaining gap.}
\textbf{UNRESOLVED}. Define the Nordheim collision functional at the atom and control its atom--density flux in an actual weak continuation.

\section{Completion Estimate}
\noindent\textbf{COMPLETION: 0\%}
This is an elementary auxiliary calculation or an explicit test of a restricted strategy. No new portion of the full source target is established.

\subsection{Sources}
\begin{itemize}
\item[\hypertarget{source:statement:1}{S1}] ULAM.AI and the UnsolvedMath Contributors, UnsolvedMath, catalogue entry. Catalogue curation under CC BY 4.0; definitions and English exposition added locally.
\url{https://huggingface.co/datasets/ulamai/UnsolvedMath}.
\item[\hypertarget{source:statement:2}{S2}] Classical and Quantum Mechanical Models of Many-Particle Systems (2023), Oberwolfach Reports, DOI 10.4171/owr/2023/38, pp. 2179--2246. Original report and mathematical context.
\url{https://doi.org/10.4171/owr/2023/38}.
\end{itemize}
\label{end:statement}
\end{document}
